\documentclass[10pt,a4paper]{amsart}
\usepackage[margin=19mm]{geometry}
\usepackage{amsmath,amssymb,mathtools}
\usepackage[scr=rsfs,cal=euler]{mathalfa}
\usepackage{xfrac}
\usepackage[unicode,hidelinks,bookmarks=false]{hyperref}
\newcommand{\ie}{i.e.\ }
\newcommand{\cf}{cf.\ }
\newcommand{\resp}{respectively\ }
\newcommand{\Subsection}{\subsection}
\providecommand{\nobreakdash}{\nobreak\hbox{-}}
\setlength{\emergencystretch}{2em}
\multlinegap0pt
\usepackage{bm}
\usepackage{bbm}
\usepackage{dsfont}
\usepackage{xspace}
\usepackage{cancel}
\usepackage{mathtools}
\usepackage{amsthm}
\makeatletter
\@ifundefined{dummy}{\newtheorem{dummy}{Dummy}}{}
\@ifundefined{rmk}{\newtheorem{rmk}[dummy]{Remark}}{}
\@ifundefined{thm}{\newtheorem{thm}[dummy]{Theorem}}{}
\makeatother
\DeclareMathOperator{\rk}{rank}
\newcommand{\consta}{{\frac{1}{2}}}
\newcommand{\oPMb}{{\overline{\mathcal{PM}}}}
\title[Open $r$-spin theory in genus one, and the Gelfand-Dikii wav]{Open $r$-spin theory in genus one, and the Gelfand-Dikii wave function}
\author{Ran J. Tessler, Yizhen Zhao}
\date{Source: arXiv:2601.04114v1 (CC BY 4.0)}
\begin{document}
\maketitle

\noindent\emph{Source and licence.} Open $r$-spin theory in genus one, and the Gelfand-Dikii wave function. Version arXiv:2601.04114v1 (CC BY 4.0), distributed under the Creative Commons Attribution 4.0 International licence, as recorded in the arXiv licence metadata for this exact version and shown on the abstract page https://arxiv.org/abs/2601.04114v1. Full rights notice: licenses/arxiv-2601.04114-CC-BY-4.0.txt.\par\smallskip

\noindent\textit{Editorial scope.} Counted unit: the Thm whose source label is \texttt{thm trr g=1 unordered} in the article (source0.tex lines 3038--3049, source label \texttt{thm trr g=1 unordered}), delivered together with the article's own complete original proof of it (source0.tex lines 3053--3120, one \texttt{proof} environment of 68 lines). No mathematics is written, repaired or re-derived: statement and proof are the article's own text line for line. The proof is detached from the statement in the source: the article states the result at lines 3038--3049 and prints its proof at lines 3053--3120, and the proof environment's own header names the statement, so the association is the article's own. Required context retained uncounted: eq rank of witten bundle g=1 (lines 998--998); eq trr closed extended (lines 1783--1788); rmk vanish without psi g=1 (lines 2393--2398); eq: bct trr 2 (lines 2511--2518); eq trr g=1 ordered (lines 2534--2542); thm indenpdent of order (lines 3019--3024); eq indep of choice g=1 (lines 3021--3023). Every definition, display and label of those lines is retained unchanged. Omitted from this package and not redistributed: 1--997, 999--1782, 1789--2392, 2399--2510, 2519--2533, 2543--3018, 3024--3037, 3050--3052, 3121--3445. Nothing inside a retained excerpt is omitted or altered: every retained body is delivered exactly as the article's lines are written, so no omission receipt of any kind accompanies this package. Citations: the retained excerpts carry no citation command at all, so no bibliography entry of the article is transcribed and no reference is redistributed. Reference adaptation: none was needed and none was made; every reference inside the delivered unit resolves inside it, so no unresolved reference or citation is produced. Printed slips kept literal: no printed word, symbol, hypothesis or number of the counted statement or of its proof was altered. Layout and class: the article's own class is \texttt{article} with packages including inputenc, geometry, amsmath, amsthm, amssymb, amscd, dsfont, bm, epic, accents, graphicx, xy, tikz, multicol; the wrapper is the lane's pinned \texttt{amsart} editorial wrapper at 10pt on A4, which restates the theorem-like environments the retained text needs and adds the article's own macro definitions that the retained text needs (\texttt{\textbackslash consta}, \texttt{\textbackslash oPMb}, \texttt{\textbackslash rk}), with extra wrapper packages (bm, bbm, dsfont, xspace, cancel, mathtools). The delivered numbering is the unit's own. Assets: no retained span contains a figure environment, a graphics-inclusion command, a PDF-page inclusion, a table or a TikZ picture; no image file is copied into this package, the package declares no asset dependency and no image file is redistributed with it. Licence: version arXiv:2601.04114v1 of this article is distributed under the Creative Commons Attribution 4.0 International licence, as recorded in the arXiv licence metadata for this exact version and shown on the abstract page https://arxiv.org/abs/2601.04114v1. A full rights notice is delivered at licenses/arxiv-2601.04114-CC-BY-4.0.txt. Characters: every retained excerpt is pure ASCII, so the delivered engine, XeLaTeX with its default Latin Modern text font, reports no missing character.
\begin{equation}\label{eq rank of witten bundle g=1}\frac{2 \sum_{i\in I} a_i + \sum_{j\in B} b_j }{r}.\end{equation}

\begin{equation}\label{eq trr closed extended}
    \left\langle\tau^{a_{i_1}}_{d_{i_1}+1}\prod_{i\in I\setminus \{i_1\}}\tau^{a_i}_{d_i}\right\rangle^{\frac{1}{r},\text{ext}}_0=
    \sum_{\substack{R_1\sqcup R_2=I\setminus\{i_1\}\\j_1,j_2\in R_2\\-1 \le a \le r-2}} 
    \left\langle\tau^a_0 \tau^{a_{i_1}}_{d_{i_1}} \prod_{i\in R_1}\tau^{a_i}_{d_i}\right\rangle^{\frac{1}{r},\text{ext}}_0 
    \left\langle\tau^{r-2-a}_0\prod_{i\in R_2}\tau^{a_i}_{d_i}\right\rangle^{\frac{1}{r},\text{ext}}_0.
\end{equation}

\begin{rmk}\label{rmk vanish without psi g=1}
    If $d_i=0$ for all $i\in I$, then the genus-one correlators are always zero because in this case by \eqref{eq rank of witten bundle g=1} we have 
    $
    \rk E=\rk \mathcal W_{1,B,I}<\dim \oPMb_{1,B,I}^{1/r}.
    $
\end{rmk}

\begin{equation}\label{eq: bct trr 2}
\begin{split}
&\left\langle \tau_{d_{i_1}+1}^{a_{i_1}}\prod_{i\in I \setminus \{i_1\}}\tau^{a_i}_{d_i}\prod_{i\in B}\sigma^{b_i}\right\rangle^{\frac{1}{r},o}_0\hspace{-0.2cm}\\=&
\sum_{\substack{R_1 \sqcup R_2 =  I\setminus \{i_1,j_I\}\\ -1\le a \le r-2}}\hspace{-0.1cm}\left\langle \tau_0^{a}\tau_{d_{i_1}}^{a_{i_1}}\prod_{i \in R_1}\tau_{d_i}^{a_i}\right\rangle^{\frac{1}{r},\text{ext}}_0
\hspace{-0.1cm}\left\langle \tau_0^{r-2-a}\tau^{a_{j_I}}_{d_{j_I}}\prod_{i\in R_2}\tau^{a_i}_{d_i}\prod_{i\in B}\sigma^{b_i}\right\rangle^{\frac{1}{r},o}_0\\
&+\hspace{-0.1cm}\sum_{\substack{R_1 \sqcup R_2 = I \setminus \{i_1,j_I\} \\ T_1 \sqcup T_2 = B}} \hspace{-0.1cm} \left\langle \tau^{a_{i_1}}_{d_{i_1}}\prod_{i \in R_1} \tau^{a_i}_{d_i}\prod_{i\in T_1}\sigma^{b_i}\right\rangle^{\frac{1}{r},o}_0 \hspace{-0.1cm}\left\langle\tau^{a_{j_I}}_{d_{j_I}} \prod_{i \in R_2} \tau^{a_i}_{d_i} \sigma^{r-2}\prod_{i\in T_2}\sigma^{b_i}\right\rangle^{\frac{1}{r}, o}_0.
\end{split}
\end{equation}

\begin{equation}\label{eq trr g=1 ordered}
\begin{split}
&\left\langle \tau_{d_{i_1}+1}^{a_{i_1}}\prod_{i\in I\setminus \{i_1\}}\tau^{a_i}_{d_i}\prod_{i\in B}\sigma^{b_i}\right\rangle^{\frac{1}{r},o,\gamma,\bm{s}}_1\hspace{-0.2cm}\\
=&\sum_{\substack{R_1 \sqcup R_2 = I\setminus\{i_1\}\\ -1\le a \le r-2}}\hspace{-0.1cm}\left\langle \tau_0^{a}\tau_{d_{i_1}}^{a_{i_1}}\prod_{i \in R_1}\tau_{d_i}^{a_i}\right\rangle^{\frac{1}{r},\text{ext}}_0
\hspace{-0.1cm}\left\langle \tau_0^{r-2-a}\prod_{i\in R_2}\tau^{a_i}_{d_i}\prod_{i\in B}\sigma^{b_i}\right\rangle^{\frac{1}{r},o,\gamma\vert_{R_2},\bm{s}\vert_{1,B,R_2\sqcup\{i^{r-2-a}_I\}}}_1\\
&+\hspace{-0.1cm}\sum_{\substack{R_1 \sqcup R_2 =  I\setminus\{i_1\} \\ T_1 \sqcup T_2 =  B}} \hspace{-0.1cm} \left\langle \tau^{a_{i_1}}_{d_{i_1}}\prod_{i \in R_1} \tau^{a_i}_{d_i}\prod_{i\in T_1}\sigma^{b_i}\right\rangle^{\frac{1}{r},o}_0 \hspace{-0.1cm}\left\langle \prod_{i \in R_2} \tau^{a_i}_{d_i} \sigma^{r-2}\prod_{i\in T_2}\sigma^{b_i}\right\rangle^{\frac{1}{r}, o,\gamma\vert_{R_2},\bm{s}\vert_{1,T_2\sqcup\{i^{r-2}_B\},R_2}}_1\\
&+\consta\left\langle \prod_{i\in I}\tau^{a_i}_{d_i}\sigma^{r-2}\prod_{i\in B}\sigma^{b_i}\right\rangle^{\frac{1}{r},o}_0,
    \end{split}
\end{equation}   

\begin{thm}\label{thm indenpdent of order}
     For any two orders $\gamma_1$ and $\gamma_2$ on the set $\{i\in I\colon d_i\ge 1\}$, and two choices of $\gamma_1$-canonical multisection $\bm{s}^1$ and $\gamma_2$-canonical multisection $\bm{s}^2$, we have
  \begin{equation}\label{eq indep of choice g=1}
  \left\langle \prod_{i\in I}\tau^{a_i}_{d_i}\prod_{i\in B}\sigma^{b_i}\right\rangle^{\frac{1}{r},o,\gamma_1,\bm{s}^1}_1=\left\langle \prod_{i\in I}\tau^{a_i}_{d_i}\prod_{i\in B}\sigma^{b_i}\right\rangle^{\frac{1}{r},o,\gamma_2,\bm{s}^2}_1.
  \end{equation}
\end{thm}

\begin{thm}[Topological Recursion Relation for $g=1$ $r$-spin theory]\label{thm trr g=1 unordered}
For any $i_1\in I$ we have a TRR with respect to $i_1$:
\begin{equation}\label{eq trr g=1 unordered}
\begin{split}
&\left\langle \tau_{d_{i_1}+1}^{a_{i_1}}\prod_{i\in I\setminus \{i_1\}}\tau^{a_i}_{d_i}\prod_{i\in B}\sigma^{b_i}\right\rangle^{\frac{1}{r},o}_1\hspace{-0.2cm}\\
=&\sum_{\substack{R_1 \sqcup R_2 = I\setminus\{i_1\}\\ -1\le a \le r-2}}\hspace{-0.1cm}\left\langle \tau_0^{a}\tau_{d_{i_1}}^{a_{i_1}}\prod_{i \in R_1}\tau_{d_i}^{a_i}\right\rangle^{\frac{1}{r},\text{ext}}_0
\hspace{-0.1cm}\left\langle \tau_0^{r-2-a}\prod_{i\in R_2}\tau^{a_i}_{d_i}\prod_{i\in B}\sigma^{b_i}\right\rangle^{\frac{1}{r},o}_1\\
&+\hspace{-0.1cm}\sum_{\substack{R_1 \sqcup R_2 =  I\setminus\{i_1\} \\ T_1 \sqcup T_2 =  B}} \hspace{-0.1cm} \left\langle \tau^{a_{i_1}}_{d_{i_1}}\prod_{i \in R_1} \tau^{a_i}_{d_i}\prod_{i\in T_1}\sigma^{b_i}\right\rangle^{\frac{1}{r},o}_0 \hspace{-0.1cm}\left\langle \prod_{i \in R_2} \tau^{a_i}_{d_i} \sigma^{r-2}\prod_{i\in T_2}\sigma^{b_i}\right\rangle^{\frac{1}{r}, o}_1\\
&+\consta\left\langle \prod_{i\in I}\tau^{a_i}_{d_i}\sigma^{r-2}\prod_{i\in B}\sigma^{b_i}\right\rangle^{\frac{1}{r},o}_0.
    \end{split}
\end{equation}
\end{thm}

\begin{proof}[Proof of Theorem \ref{thm indenpdent of order} and Theorem \ref{thm trr g=1 unordered}]
    We prove the \eqref{eq indep of choice g=1}, and hence \eqref{eq trr g=1 unordered}, by an induction on $D=\sum_{i\in I} d_i$.

    The case $D=0$ is trivially true: according to Remark \ref{rmk vanish without psi g=1} the correlators in this case are always zero.

    In the case $D=1$, there is only one element $i_1\in I$ with $d_{i_1}=1>0$. Note that in this case $I^{d>0}=\{i_1\}$ is a single-element set, so the choice of the order $\gamma$ is unique. We apply \eqref{eq trr g=1 ordered}  to the correlator $\left\langle \tau_{1}^{a_{i_1}}\prod_{i\in I\setminus \{i_1\}}\tau^{a_i}_{0}\prod_{i\in B}\sigma^{b_i}\right\rangle^{\frac{1}{r},o,\gamma,\bm{s}}_1$ and express it as 
    \begin{equation}\label{eq apply trr D=1}
        \begin{split}
    &\sum_{\substack{R_1 \sqcup R_2 = I\setminus\{i_1\}\\ -1\le a \le r-2}}\hspace{-0.1cm}\left\langle \tau_0^{a}\tau_{0}^{a_{i_1}}\prod_{i \in R_1}\tau_{0}^{a_i}\right\rangle^{\frac{1}{r},\text{ext}}_0
    \hspace{-0.1cm}\left\langle \tau_0^{r-2-a}\prod_{i\in R_2}\tau^{a_i}_{0}\prod_{i\in B}\sigma^{b_i}\right\rangle^{\frac{1}{r},o,\gamma\vert_{R_2},\bm{s}\vert_{R_2\sqcup\{i^{r-2-a}_I\},B}}_1\\
    &+\hspace{-0.1cm}\sum_{\substack{R_1 \sqcup R_2 =  I\setminus\{i_1\} \\ T_1 \sqcup T_2 =  B}} \hspace{-0.1cm} \left\langle \tau^{a_{i_1}}_{0}\prod_{i \in R_1} \tau^{a_i}_{0}\prod_{i\in T_1}\sigma^{b_i}\right\rangle^{\frac{1}{r},o}_0 \hspace{-0.1cm}\left\langle \prod_{i \in R_2} \tau^{a_i}_{0} \sigma^{r-2}\prod_{i\in T_2}\sigma^{b_i}\right\rangle^{\frac{1}{r}, o,\gamma\vert_{R_2},\bm{s}\vert_{R_2,B\sqcup\{i^{r-2}_B\}}}_1\\
    &+\consta\left\langle \prod_{i\in I}\tau^{a_i}_{0}\sigma^{r-2}\prod_{i\in B}\sigma^{b_i}\right\rangle^{\frac{1}{r},o}_0.
    \end{split}
    \end{equation}
    By Remark \ref{rmk vanish without psi g=1} all the genus-one factors in \eqref{eq apply trr D=1} vanishes, the only term left is 
    $$
    \consta\left\langle \prod_{i\in I}\tau^{a_i}_{0}\sigma^{r-2}\prod_{i\in B}\sigma^{b_i}\right\rangle^{\frac{1}{r},o}_0
    $$
    which is independent of the choice of $\bm{s}$.


    Assuming we have proven \eqref{eq indep of choice g=1} and  \eqref{eq trr g=1 unordered} in in all the cases where $\sum_i d_i\le D_0$ for a integer $D_0\ge 1$,  we now prove \eqref{eq indep of choice g=1} and  \eqref{eq trr g=1 unordered} in the case $\sum_i d_i= D_0+1$. We take $i_1=\gamma_1(I)$ and use \eqref{eq trr g=1 ordered}  to the correlator 
    $\left\langle \tau_{d_{i_1}}^{a_{i_1}}\prod_{i\in I\setminus \{i_1\}}\tau^{a_i}_{d_i}\prod_{i\in B}\sigma^{b_i}\right\rangle^{\frac{1}{r},o,\gamma_1,\bm{s}^1}_1$ on the left-hand side of \eqref{eq indep of choice g=1}  as
    \begin{equation}\label{eq apply trr D>1}
        \begin{split}
    &\sum_{\substack{R_1 \sqcup R_2 = I\setminus\{i_1\}\\ -1\le a \le r-2}}\hspace{-0.1cm}\left\langle \tau_0^{a}\tau_{d_{i_1}-1}^{a_{i_1}}\prod_{i \in R_1}\tau_{d_i}^{a_i}\right\rangle^{\frac{1}{r},\text{ext}}_0
    \hspace{-0.1cm}\left\langle \tau_0^{r-2-a}\prod_{i\in R_2}\tau^{a_i}_{d_i}\prod_{i\in B}\sigma^{b_i}\right\rangle^{\frac{1}{r},o,\gamma_1\vert_{R_2},\bm{s}^1\vert_{R_2\sqcup\{i^{r-2-a}_I\},B}}_1\\
    &+\hspace{-0.1cm}\sum_{\substack{R_1 \sqcup R_2 =  I\setminus\{i_1\} \\ T_1 \sqcup T_2 =  B}} \hspace{-0.1cm} \left\langle \tau^{a_{i_1}}_{d_{i_1}-1}\prod_{i \in R_1} \tau^{a_i}_{d_i}\prod_{i\in T_1}\sigma^{b_i}\right\rangle^{\frac{1}{r},o}_0 \hspace{-0.1cm}\left\langle \prod_{i \in R_2} \tau^{a_i}_{d_i} \sigma^{r-2}\prod_{i\in T_2}\sigma^{b_i}\right\rangle^{\frac{1}{r}, o,\gamma_1\vert_{R_2},\bm{s}^1\vert_{R_2,B\sqcup\{i^{r-2}_B\}}}_1\\
    &+\consta\left\langle \tau_{d_{i_1}-1}^{a_1}\prod_{i\in I\setminus\{i_1\}}\tau^{a_i}_{d_i}\sigma^{r-2}\prod_{i\in B}\sigma^{b_i}\right\rangle^{\frac{1}{r},o}_0.
    \end{split}
    \end{equation}
    Note that every genus-one factor in \eqref{eq apply trr D>1} is independent of any choice of the order or the canonical multisection by inductive hypothesis as $\sum_{i\in R_2}d_i \le D_0$, therefore we can rewrite \eqref{eq apply trr D>1} as 
    \begin{equation}\label{eq apply trr D>1 v1}
        \begin{split}
    &\sum_{\substack{R_1 \sqcup R_2 = I\setminus\{i_1\}\\ -1\le a \le r-2}}\hspace{-0.1cm}\left\langle \tau_0^{a}\tau_{d_{i_1}-1}^{a_{i_1}}\prod_{i \in R_1}\tau_{d_i}^{a_i}\right\rangle^{\frac{1}{r},\text{ext}}_0
    \hspace{-0.1cm}\left\langle \tau_0^{r-2-a}\prod_{i\in R_2}\tau^{a_i}_{d_i}\prod_{i\in B}\sigma^{b_i}\right\rangle^{\frac{1}{r},o}_1\\
    &+\hspace{-0.1cm}\sum_{\substack{R_1 \sqcup R_2 =  I\setminus\{i_1\} \\ T_1 \sqcup T_2 =  B}} \hspace{-0.1cm} \left\langle \tau^{a_{i_1}}_{d_{i_1}-1}\prod_{i \in R_1} \tau^{a_i}_{d_i}\prod_{i\in T_1}\sigma^{b_i}\right\rangle^{\frac{1}{r},o}_0 \hspace{-0.1cm}\left\langle \prod_{i \in R_2} \tau^{a_i}_{d_i} \sigma^{r-2}\prod_{i\in T_2}\sigma^{b_i}\right\rangle^{\frac{1}{r}, o}_1\\
    &+\consta\left\langle \tau_{d_{i_1}-1}^{a_1}\prod_{i\in I\setminus\{i_1\}}\tau^{a_i}_{d_i}\sigma^{r-2}\prod_{i\in B}\sigma^{b_i}\right\rangle^{\frac{1}{r},o}_0.
    \end{split}
    \end{equation} 
    This implies that the correlator 
    $\left\langle \tau_{d_{i_1}}^{a_{i_1}}\prod_{i\in I\setminus \{i_1\}}\tau^{a_i}_{d_i}\prod_{i\in B}\sigma^{b_i}\right\rangle^{\frac{1}{r},o,\gamma_1,\bm{s}^1}_1$ does not depend on the choice of $\bm{s}^1$. Note that for the moment it still depends on $\gamma_1$ since $i_1$ in \eqref{eq apply trr D>1 v1} is chosen as $i_1=\gamma_1(I)$.

    We can also take $i_2:=\gamma_2(I)$ and express the correlator 
    $\left\langle \tau_{d_{i_2}}^{a_{i_2}}\prod_{i\in I\setminus \{i_2\}}\tau^{a_i}_{d_i}\prod_{i\in B}\sigma^{b_i}\right\rangle^{\frac{1}{r},o,\gamma_2,\bm{s}^2}_1$ on the right-hand side of \eqref{eq indep of choice g=1} using \eqref{eq trr g=1 ordered} as
    \begin{equation}\label{eq apply trr D>1 v2}
        \begin{split}
    &\sum_{\substack{R_1 \sqcup R_2 = I\setminus\{i_2\}\\ -1\le a \le r-2}}\hspace{-0.1cm}\left\langle \tau_0^{a}\tau_{d_{i_2}-1}^{a_{i_2}}\prod_{i \in R_1}\tau_{d_i}^{a_i}\right\rangle^{\frac{1}{r},\text{ext}}_0
    \hspace{-0.1cm}\left\langle \tau_0^{r-2-a}\prod_{i\in R_2}\tau^{a_i}_{d_i}\prod_{i\in B}\sigma^{b_i}\right\rangle^{\frac{1}{r},o}_1\\
    &+\hspace{-0.1cm}\sum_{\substack{R_1 \sqcup R_2 =  I\setminus\{i_2\} \\ T_1 \sqcup T_2 =  B}} \hspace{-0.1cm} \left\langle \tau^{a_{i_2}}_{d_{i_2}-1}\prod_{i \in R_1} \tau^{a_i}_{d_i}\prod_{i\in T_1}\sigma^{b_i}\right\rangle^{\frac{1}{r},o}_0 \hspace{-0.1cm}\left\langle \prod_{i \in R_2} \tau^{a_i}_{d_i} \sigma^{r-2}\prod_{i\in T_2}\sigma^{b_i}\right\rangle^{\frac{1}{r}, o}_1\\
    &+\consta\left\langle \tau^{a_{i_2}}_{d_{i_2}-1}\prod_{i\in I\setminus \{i_2\}}\tau^{a_i}_{d_i}\sigma^{r-2}\prod_{i\in B}\sigma^{b_i}\right\rangle^{\frac{1}{r},o}_0,
    \end{split}
    \end{equation}
    which only depends on $i_2=\gamma_2(I)$. Therefore \eqref{eq indep of choice g=1} is proven in the case $\gamma_1(I)=\gamma_2(I)$.

    To conclude the proof, we need to show that \eqref{eq apply trr D>1 v1} coincides with \eqref{eq apply trr D>1 v2} even if $i_1=\gamma_1(I)\ne \gamma_2(I) = i_2$. This can be done by applying the Topological Recursion Relations \eqref{eq trr closed extended}, \eqref{eq: bct trr 2} and \eqref{eq trr g=1 unordered} to the closed extended, open genus-zero, and open genus-one correlators in  \eqref{eq apply trr D>1 v1} and \eqref{eq apply trr D>1 v2}, in the following way. 
    \begin{itemize}
        \item[-] For the closed extended correlators $\left\langle \tau_0^{a}\tau_{d_{i_1}-1}^{a_{i_1}}\prod_{i \in R_1}\tau_{d_i}^{a_i}\right\rangle^{1/r,\text{ext}}_0$ in \eqref{eq apply trr D>1 v1}, if $i_2\notin R_1$, we leave it unchanged; if $i_2\in R_1$, we apply \eqref{eq trr closed extended} with respect to $(i_2,\{i_1,i_I^a\})$, where $i_I^a$ is the internal tail corresponding to $\tau_0^a$ in the correlator.
        \item[-] For the genus-zero open correlators of the form $\left\langle \tau_{d_{i_1}-1}^{a_{i_1}}\prod_{i \in I'}\tau_{d_i}^{a_i} \prod_{i\in B'}\sigma^{b_i}\right\rangle^{1/r,o}_0$ in \eqref{eq apply trr D>1 v1}, if $i_2\notin I'$, we leave it unchanged; if $i_2\in I'$, we apply \eqref{eq: bct trr 2} with respect to $(i_2,i_1)$.
        \item[-] For the genus-one open correlators of the form $\left\langle \prod_{i \in I'}\tau_{d_i}^{a_i} \prod_{i\in B'}\sigma^{b_i}\right\rangle^{1/r,o}_1$ in \eqref{eq apply trr D>1 v1}, if $i_2\notin I'$, we leave it unchanged; if $i_2\in I'$, we apply \eqref{eq trr g=1 unordered} with respect to $i_2$. Note that \eqref{eq trr g=1 unordered} is already proven in this case by inductive hypothesis. 
    \end{itemize}

    After substituting these correlators with the results after applying TRRs, we get a new expression for \eqref{eq apply trr D>1 v1}. Note that in this expression $i_1$ and $i_2$ are symmetry.

    We can also apply TRRs to the correlators in \eqref{eq apply trr D>1 v2} in the same way after exchanging the roles of $i_1$ and $i_2$. In this way we get a new expression for \eqref{eq apply trr D>1 v2}, which is exactly the same as the new expression for \eqref{eq apply trr D>1 v1}. This proves \eqref{eq indep of choice g=1}.

    
\end{proof}

\end{document}
